% !TeX root = ../thesis.tex
\chapter{Abkürzungsverzeichnis}

% Alphabetisch sortiert; neue Abkürzungen hier definieren und im Text mit \ac markieren.
\begin{acronym}[SHA-256]
  \acro{anls}[ANLS]{Average Normalized Levenshtein Similarity}
  \acro{api}[API]{Application Programming Interface}
  \acro{bf16}[BF16]{Brain Floating Point (16-Bit)}
  \acro{cer}[CER]{Character Error Rate}
  \acro{cli}[CLI]{Command-Line Interface}
  \acro{cpu}[CPU]{Central Processing Unit}
  \acro{csp}[CSP]{Content Security Policy}
  \acro{csv}[CSV]{Comma-Separated Values}
  \acro{ctc}[CTC]{Connectionist Temporal Classification}
  \acro{cuda}[CUDA]{Compute Unified Device Architecture}
  \acro{dia}[DIA]{Document Image Analysis}
  \acro{dla}[DLA]{Document Layout Analysis}
  \acro{dpi}[DPI]{Dots Per Inch}
  \acro{ece}[ECE]{Expected Calibration Error}
  \acro{erp}[ERP]{Enterprise Resource Planning}
  \acro{gnhk}[GNHK]{GoodNotes Handwriting Kollection}
  \acro{gpu}[GPU]{Graphics Processing Unit}
  \acro{html}[HTML]{HyperText Markup Language}
  \acro{htr}[HTR]{Handwritten Text Recognition}
  \acro{id}[ID]{Identifikator}
  \acro{json}[JSON]{JavaScript Object Notation}
  \acro{ki}[KI]{Künstliche Intelligenz}
  \acro{llm}[LLM]{Large Language Model}
  \acro{nfc}[NFC]{Normalization Form Canonical Composition}
  \acro{ocr}[OCR]{Optical Character Recognition}
  \acro{oom}[OOM]{Out of Memory}
  \acro{osd}[OSD]{Orientation and Script Detection}
  \acro{pdf}[PDF]{Portable Document Format}
  \acro{psm}[PSM]{Page Segmentation Mode}
  \acro{ram}[RAM]{Random-Access Memory}
  \acro{sdpa}[SDPA]{Scaled Dot-Product Attention}
  \acro{sha256}[SHA-256]{Secure Hash Algorithm (256-Bit)}
  \acro{trocr}[TrOCR]{Transformer-based Optical Character Recognition}
  \acro{vlm}[VLM]{Vision-Language Model}
  \acro{vram}[VRAM]{Video Random-Access Memory}
  \acro{wer}[WER]{Word Error Rate}
  \acro{xlsx}[XLSX]{Office Open XML Spreadsheet}
\end{acronym}
